\section{相关工作}

\textbf{2D目标检测（2D Object Detection）。} RCNN~\cite{girshick2014rcnn}开创了使用深度学习进行目标检测的先河。它将一组预选的目标候选框输入卷积神经网络（CNN），并据此预测边界框参数。尽管该方法表现出令人惊讶的性能，但由于需要对每个目标候选框执行一次卷积网络前向传播，其速度比其他方法慢一个数量级。为解决此问题，Fast RCNN~\cite{ren2015faster}引入了一个共享的可学习CNN，在单次前向传播中处理整张图像。为进一步提升性能和速度，Faster RCNN~\cite{ren2015faster}包含了一个区域提议网络（Region Proposal Network, RPN），该网络与检测网络共享全图卷积特征，从而几乎零成本地生成区域提议。Mask RCNN~\cite{he2017maskrcnn}引入了一个掩码预测分支，以实现并行实例分割。这些方法通常涉及多阶段优化，在实践中可能速度较慢。与这些多阶段方法不同，SSD~\cite{liu2016ssd}和YOLO~\cite{Redmon2015YouOL}以单次方式执行密集预测。虽然它们显著快于上述替代方案，但仍依赖NMS去除冗余框预测。这些方法相对于预定义的锚框（Anchors）预测边界框。CenterNet~\cite{zhou2019centernet}和FCOS~\cite{tian2019fcos}改变了范式，从逐锚框预测转向逐像素预测，大大简化了常见的目标检测流程。

\textbf{基于集合的目标检测（Set-based Object Detection）。} DETR~\cite{detr}将目标检测视为一个集合到集合问题。它使用Transformer~\cite{Vaswani2017}捕获特征和目标交互。DETR学习将预测结果分配给一组真值框，因此不需要后处理来过滤冗余框。然而，DETR的一个关键缺点是需要大量的训练时间。Deformable DETR~\cite{zhu2021deformable}分析了DETR收敛缓慢的原因，并提出了一种可变形自注意力模块来定位特征并加速训练。与此同时，\cite{Sun2020RethinkingTS}将DETR收敛缓慢归因于集合损失和Transformer交叉注意力机制，提出了TSP-FCOS和TSP-RCNN两种变体来克服这些问题。SparseRCNN~\cite{peize2020sparse}将集合预测融入RCNN风格的流程，性能优于多阶段目标检测且无需NMS。OneNet~\cite{peize2020onenet}研究了一个有趣的现象：密集目标检测器配备最小代价集合损失后，可实现无NMS检测。在3D领域，Object DGCNN~\cite{objdgcnn}研究了从点云进行3D目标检测。它将3D目标检测建模为动态图上的消息传递，将DGCNN框架推广到预测一组目标。与DETR类似，Object DGCNN也是无NMS的。

\textbf{单目3D目标检测（Monocular 3D Object Detection）。} 从RGB图像进行3D检测的早期方法是Mono3D~\cite{chen2016monocular}，它利用语义和形状线索，在训练时使用场景约束和额外先验从一组3D候选框中进行选择。\cite{roddick2018orthographic}使用鸟瞰视图（Bird's-Eye View, BEV）进行单目3D检测，而~\cite{mousavian20173d}利用2D检测结果，通过最小化2D-3D投影误差进行3D边界框回归。使用2D检测器作为3D计算的起点最近已成为一种标准方法~\cite{kehl2017ssd,ku2019monocular}。其他工作还探索了可微渲染~\cite{beker2020monocular}或3D关键点检测~\cite{barabanau2019monocular,liu2020smoke,zhou2019centernet}等方面的进展，以实现当前最优的3D目标检测性能。所有这些方法都是在单目设置下运行的，扩展到多相机的方式是独立处理每个视角的图像，然后在后处理阶段合并输出。